\section{Choosing the smoothing parameters - Cross Validation}
\label{sec:auto}

The smoothing parameter
$\balpha$
can be chosen automatically~\cite{Fessler:91:lin}
by minimizing the 
cross validation (CV) score:
\[
\balpha_{\CV} \defequ \arg \min_{\balpha} \CV(\balpha),
\]
\be
\CV(\balpha) \defequ \frac{1}{N} \sum_{n=1}^N 
(\meas_n - \mfun(\hat{\g}_{\balpha, -n}(t_n)))' \mcov_n^{-1}
(\meas_n - \mfun(\hat{\g}_{\balpha, -n}(t_n)))
.\ee{eq:cv-score}
$\hat{\g}_{\balpha, -n}$
is the solution to the smoothing problem (\ref{eq:multi-min})
with N-1 data points,
posed
while excluding
the pair $(t_n, \meas_n)$.
Each data pair is dropped in turn,
the smoothed curve $\bghatan$ is estimated,
and the predicted value $\bghatan(t_n)$
is compared with the unused measurement $\y_n$.
If the CV score is small,
then we have
chosen the smoothing parameter
that makes
the estimated curve a good self predictor.

Although equation~(\ref{eq:cv-score})
illustrates the idea behind cross validation,
it is computationally inefficient.
We have shown that,
in the linear case,
one can rewrite (\ref{eq:cv-score})
as:
\be
\CV(\balpha) = \frac{1}{N} \sum_{n=1}^N
\| \mcov_n^{-\frac{1}{2}} (\I_M - \Ad{nn})\inv
	(\meas_n - \bghata(t_n)) \|^2
,\ee{eq:cv-score2}
where $\Ad{nn}$ is the $n$'th $M \times M$
block diagonal submatrix of the influence matrix.
In the nonlinear case,
(\ref{eq:cv-score2}) is a useful approximation,
as shown in~\cite{Fessler:91:non}.
The properties of this approximation
are not completely understood.

By using the Hutchinson and de Hoog algorithm~\cite{HutchHoog:85},
(\ref{eq:cv-score2}) is computed
in only O($M^3N$) operations.

Subroutine
{\tt bestcv},
described below,
computes
$\balpha_{\CV}$
and
$\hat{\g}_{\balpha_{\CV}}$.
The VSPLINE library also
contains routines for computing
Generalized Cross Validation
and Unbiased Risk scores~\cite{Fessler:91:lin}.
